\section{Personas}

\subsection{Marco}
\begin{table}[h]
\centering
\begin{tabular}{|l|p{10cm}|}
\hline
\textbf{Età} & 23 anni \\
\hline
\textbf{Occupazione} & Studente universitario del corso di laurea magistrale in Ingegneria e Scienze Informatiche \\
\hline
\textbf{Titolo di studi} & Laurea triennale \\
\hline
\end{tabular}
\end{table}
Marco è uno studente che frequenta regolarmente le lezioni e utilizza quotidianamente il computer per studiare. Durante le lezioni prende appunti digitali, ma spesso li trova disorganizzati e difficili da collegare alle slide del corso.

\subsubsection{Obiettivi}
\begin{itemize}
    \item Prendere appunti in modo ordinato e strutturato
    \item Associare direttamente gli appunti alle slide del corso
    \item Accedere al materiale di studio da diversi dispositivi
    \item Generare riassunti automatici per il ripasso pre-esame
\end{itemize}

\subsubsection{Scenario d'uso}
\begin{itemize}
    \item \textbf{Motivazioni}\par
    Marco è motivato dalla necessità di prendere appunti in modo ordinato e facilmente consultabile, riducendo il tempo necessario alla riorganizzazione e al ripasso del materiale in vista degli esami.
    \item \textbf{Contesto}\par
    Lo scenario si svolge durante una lezione universitaria in aula. Marco utilizza il proprio laptop collegato alla rete Wi-Fi dell'università e accede alla Web App Easy Notes tramite browser.
    \item \textbf{Azioni}\par
    \begin{enumerate}
        \item Marco accede a Easy Notes
        \item Seleziona il corso a cui è iscritto e apre le slide della lezione.
        \item Per ogni slide visualizzata, inserisce appunti testuali nella sezione dedicata.
        \item Al termine della lezione salva automaticamente appunti e slide sul cloud
        \item Arrivato a casa, Marco utilizza la funzionalità di riassunto automatico per prepararsi allo studio.
    \end{enumerate}
\end{itemize}

\subsection{Serena}
\begin{table}[h]
\centering
\begin{tabular}{|l|p{10cm}|}
\hline
\textbf{Età} & 26 anni \\
\hline
\textbf{Occupazione} & Studentessa lavoratrice iscritta al corso di laurea in Scienze e Tecnologie Psicologiche. \\
\hline
\textbf{Titolo di studi} & Diploma di scuola superiore \\
\hline
\end{tabular}
\end{table}

Serena frequenta l'università mentre lavora part-time. Non riesce sempre a seguire le lezioni in presenza e studia prevalentemente da casa, utilizzando slide e materiale caricato dai docenti.

\subsubsection{Obiettivi}
\begin{itemize}
    \item Studiare in modo efficiente con tempi ridotti
    \item Avere appunti sempre sincronizzati e accessibili
    \item Riassumere rapidamente grandi quantità di contenuto
\end{itemize}

\subsubsection{Scenario d'uso}
\begin{itemize}
    \item \textbf{Motivazioni}\par
    Serena è motivata dall'esigenza di ottimizzare il tempo a disposizione per lo studio e di mantenere il materiale sempre aggiornato, nonostante l'impossibilità di seguire tutte le lezioni in presenza.
    \item \textbf{Contesto}\par
    Lo scenario si svolge a casa, in un momento serale. Serena utilizza il proprio laptop personale e accede a Easy Notes tramite browser, con connessione Internet domestica.
    \item \textbf{Azioni}\par
    \begin{enumerate}
        \item Serena accede a Easy Notes.
        \item Riceve una notifica da Easy Notes relativa al caricamento di nuove slide da parte del docente per la lezione di domani. 
        \item Apre il materiale del corso e aggiunge appunti personali associandoli alle slide. 
        \item Successivamente utilizza la funzione di riassunto automatico per ottenere una sintesi del contenuto e pianificare lo studio nei giorni successivi.
    \end{enumerate}
\end{itemize}

\subsection{Prof. Bianchi}
\begin{table}[h]
\centering
\begin{tabular}{|l|p{10cm}|}
\hline
\textbf{Età} & 48 anni \\
\hline
\textbf{Occupazione} & Docente universitario \\
\hline
\textbf{Titolo di studi} & Dottorato di ricerca \\
\hline
\end{tabular}
\end{table}

Il professor Bianchi tiene corsi universitari con numerosi studenti. Utilizza slide digitali per le lezioni e desidera una piattaforma semplice per distribuire il materiale didattico e mantenere una comunicazione efficace con gli studenti.

\subsubsection{Obiettivi}
\begin{itemize}
    \item Caricare e gestire le slide del corso
    \item Organizzare il materiale in modo chiaro
    \item Notificare gli studenti in caso di nuovi contenuti
    \item Ridurre la dispersione delle informazioni
\end{itemize}

\subsubsection{Scenario d'uso}
\begin{itemize}
    \item \textbf{Motivazioni}\par
    Il professor Bianchi è motivato dal desiderio di distribuire in modo efficace il materiale didattico e di mantenere gli studenti costantemente aggiornati sui contenuti del corso.
    \item \textbf{Contesto}\par
    Lo scenario si svolge nel suo ufficio universitario. Il docente utilizza un computer desktop con connessione Internet e accede alla Web App Easy Notes tramite browser.
    \item \textbf{Azioni}\par
    \begin{enumerate}
        \item Il professore accede a Easy Notes
        \item Seleziona la classe di riferimento e carica le nuove slide della lezione.
        \item Una volta completato il caricamento, il sistema invia automaticamente una notifica agli studenti iscritti al corso. 
        \item Il docente verifica che il materiale sia correttamente disponibile sulla piattaforma.
    \end{enumerate}
\end{itemize}